% Compile with XeLaTeX (uses the system Times New Roman font).
\documentclass[12pt]{article}

\usepackage[a4paper,margin=2.3cm]{geometry}
\usepackage{amsmath,amssymb,amsthm}
\usepackage{fontspec}
\setmainfont{Times New Roman}
\usepackage{graphicx}
\usepackage{booktabs}
\usepackage{xcolor}
\usepackage[hidelinks]{hyperref}
\usepackage[font=small,skip=3pt]{caption}

\graphicspath{{data/output/}{../data/output/}}
\setlength{\parskip}{2pt}

\DeclareMathOperator{\tr}{tr}
\newcommand{\R}{\mathbb{R}}
\newcommand{\Xt}{\widetilde{X}}
\newcommand{\xt}{\widetilde{x}}
\newcommand{\Nt}{\widetilde{N}}
\newcommand{\norm}[1]{\lVert #1 \rVert}
\newcommand{\T}{^{\mathsf{T}}}

\theoremstyle{definition}\newtheorem{definition}{Definition}
\theoremstyle{plain}\newtheorem{proposition}{Proposition}

% ---- compact title block (no separate title page) ----
\title{\vspace{-1.2cm}\bfseries Question 3 --- Design, Covariance and Gram Matrices:\\[2pt]
\large Cleaning, Principal Component Analysis, and Clustering}
\author{SIMC 2026 Endeavour Challenge \quad(Parts a--e)}
\date{}

\begin{document}
\maketitle
\vspace{-0.7cm}

\begin{abstract}\noindent\small
We analyse a design matrix of $N=1000$ infrared spectra at $M=3401$ frequencies. Thresholding
the row-norm distribution at its spectral gaps removes $40$ empty vials and $36$ high-norm
outliers, leaving $\Nt=924$ spectra. The covariance matrix $C=\tfrac1{\Nt}\Xt\T\Xt$ has total
variance $\tr C\approx7.90$, dominance ratio $\lambda_1/\lambda_2\approx3.53$, and a sharp
elbow at $k^\star=8$ signal eigenvalues; the ``$95\%$ of variance'' rule is shown to be
misleading (it needs $667$ components). The leading eigenvector is an interpretable
difference-spectrum, and projection onto the top four eigenvectors resolves the data into
$\approx25$ chemical concoctions.
\end{abstract}

\section{Setup}\label{sec:setup}
$X\in\R^{N\times M}$ has rows = \emph{samples} (vials) and columns = \emph{features}
(frequencies); $X_{ij}$ is the absorbance of sample $i$ at frequency $j$. The supplied array
has shape $3401\times1000$: it is stored \emph{transposed}, and its feature axis is $3401$, not
the $3041$ quoted in the prompt (a digit transposition). We transpose so rows index samples and
use $M=3401$. The \emph{Euclidean norm} $\norm{x}=(\sum_j x_j^2)^{1/2}$ summarises a spectrum's
overall magnitude.

\section{Part (a): cleaning}\label{sec:clean}
Empty vials record only background noise: small, flat, hence small norm. Sorting the $1000$ norms
reveals a clean gap --- $40$ vials at $\norm{x}\approx1.6$, then a void until the bulk at
$\norm{x}\approx4.3$ --- so we threshold at the gap midpoint (a robust, constant-free choice).
This flags exactly $T=40$ empties (rows $960$--$999$). The upper tail holds a \emph{second} error:
$36$ vials at $\norm{x}\approx8.9$ (rows $924$--$959$), detached by a gap $\approx5\times$ wider
than any internal gap. We retain $i$ iff $\tau_{\mathrm{lo}}\!\le\!\norm{x_i}\!\le\!\tau_{\mathrm{hi}}$
via a Boolean keep-mask (non-destructive, composable, index-preserving), leaving $\Nt=924$ spectra.
Validation: removed spectra are flat while retained ones are peaked, and empties are $50.2\%$ exact
zeros versus $24.5\%$ (background noise clipped at zero is exactly $0$ with probability $\approx\!\tfrac12$).
Section~\ref{sec:clusters} shows the $36$ outliers form one isolated cluster whose removal drops the
cluster count by exactly one, so the decision is clean (and reversible).

\begin{figure}[h]\centering
\includegraphics[width=0.80\linewidth]{q3a_part_i.png}
\caption{Row-norm distribution: $40$ empty vials below the gap; example empty (flat) vs genuine
(peaked) spectra.}\label{fig:clean}
\end{figure}

\section{Part (b): variance and the covariance matrix}\label{sec:cov}
Mean-centering, $\Xt_{ij}=X_{ij}-\bar X_j$ with $\bar X_j=\tfrac1{\Nt}\sum_i X_{ij}$, removes the
shared mean spectrum so the analysis describes how samples \emph{differ}. The covariance matrix is
$C=\tfrac1{\Nt}\Xt\T\Xt\in\R^{M\times M}$.

\begin{proposition}\label{prop:diag}
$C$ is symmetric and $C_{jj}=\tfrac1{\Nt}\sum_k\Xt_{kj}^2=\mathrm{Var}_j$, the variance of feature $j$.
\end{proposition}
\begin{proof}
$C\T=\tfrac1{\Nt}(\Xt\T\Xt)\T=C$; setting $i=j$ and using zero column-means gives the variance.
\end{proof}

\begin{proposition}\label{prop:trace}
$\tr C=\sum_j\mathrm{Var}_j$ (total variance) is invariant under orthogonal change of basis, and
equals $\sum_i\lambda_i$.
\end{proposition}
\begin{proof}
$\tr(Q\T C Q)=\tr(CQQ\T)=\tr C$ for $Q\T Q=I$; diagonalising $C=Q\Lambda Q\T$ gives $\tr C=\sum_i\lambda_i$.
\end{proof}

Empirically $\tr C\approx7.90$. Variance concentrates in a few narrow molecular bands
(Figure~\ref{fig:var}); the maximum is at $j^\star=910$ with $C_{j^\star j^\star}\approx0.0199$.
Most features are flat, so the data is effectively low-dimensional.

\begin{figure}[h]\centering
\includegraphics[width=0.82\linewidth]{q3b_variance.png}
\caption{Per-feature variance $C_{jj}$; the maximum is at $j^\star=910$ (red).}\label{fig:var}
\end{figure}

\section{Part (c): the eigenspectrum (PCA)}\label{sec:pca}
\begin{proposition}\label{prop:rayleigh}
For unit $w$, the variance of the data along $w$ is $w\T C w$; if $C\nu=\lambda\nu$ with $\norm{\nu}=1$
then $\nu\T C\nu=\lambda$. Moreover $C$ is positive semi-definite, so every $\lambda_i\ge0$.
\end{proposition}
\begin{proof}
The projections $z_i=\xt_i\!\cdot w$ have variance $\tfrac1{\Nt}\sum_i z_i^2=w\T C w$; substituting an
eigenvector gives $\lambda$. And $w\T C w=\tfrac1{\Nt}\norm{\Xt w}^2\ge0$, so $\lambda=\nu\T C\nu\ge0$.
\end{proof}

By the spectral theorem $C$ has real eigenvalues and an orthonormal eigenbasis; we use
\texttt{numpy.linalg.eigh} (symmetric solver: real eigenvalues, orthonormal eigenvectors, faster and
stabler than \texttt{eig}), sort decreasingly, and clip rounding-level negatives to $0$.

\begin{table}[h]\centering\small
\begin{tabular}{lccccccccc}
\toprule
$k$ & 1 & 2 & 3 & 4 & 5 & 6 & 7 & 8 & 9\\
$\lambda_k$ & 2.492 & 0.705 & 0.590 & 0.254 & 0.123 & 0.105 & 0.078 & 0.039 & 0.0093\\
\bottomrule
\end{tabular}
\caption{Leading eigenvalues. The drop $\lambda_8/\lambda_9\approx4.1$ is the largest in the spectrum.}
\label{tab:eig}
\end{table}

The dominance ratio is $\lambda_1/\lambda_2\approx3.53$: one direction dominates, indicating strong
low-dimensional structure (a wide spectral gap, the condition for fast Power-Method convergence,
Question 2). The \emph{elbow} is the largest multiplicative drop between consecutive eigenvalues
(a ratio, since the signal/noise boundary is multiplicative on the log scale of
Figure~\ref{fig:scree}): $\lambda_8/\lambda_9\approx4.1$ gives $k^\star=8$. Here $\lambda_8$ remains
$\approx4\times$ above the noise floor (signal) while $\lambda_9$ lies on it (noise); this matches the
$\sim10$ dominant molecules. The eight signal components explain only $\approx55\%$ of the variance,
and reaching $95\%$ needs $k=667$ components (Table~\ref{tab:cum}), because the noise floor is spread
over $\approx900$ dimensions; the ``$95\%$'' rule therefore overstates the dimensionality, and the
elbow is the honest measure. Check: $\sum_i\lambda_i\approx\tr C\approx7.90$ (Prop.~\ref{prop:trace}).

\begin{table}[h]\centering\small
\begin{tabular}{lccccc}
\toprule
cumulative variance & $50\%$ & $80\%$ & $90\%$ & $95\%$ & $99\%$\\
components $k$ & 4 & 309 & 516 & 667 & 850\\
\bottomrule
\end{tabular}
\caption{Components to reach each threshold: contrast the $8$ signal components ($55\%$) with $667$ for $95\%$.}
\label{tab:cum}
\end{table}

\begin{figure}[h]\centering
\includegraphics[width=0.84\linewidth]{q3c_spectrum.png}
\caption{Eigenvalue spectrum (linear and log). The log plot shows a sharp elbow at $k^\star=8$,
then a flat noise floor.}\label{fig:scree}
\end{figure}

\section{Part (d): the leading eigenvector}\label{sec:lead}
$\nu_1\in\R^M$ (feature space) gives each frequency a \emph{loading}; the \emph{scores}
$z_i=\xt_i\!\cdot\nu_1$ are the maximally-spread one-dimensional summaries (Prop.~\ref{prop:rayleigh}).
Eigenvectors are sign-ambiguous; we fix the largest loading positive. The loadings are supported on
two narrow bands ($j\approx900$--$960$, $3150$--$3260$) with $\approx27\%$ negative entries, so $\nu_1$
is a \emph{contrast} (balance of absorption between bands), not overall intensity. With cosine
similarity $\cos(a,b)=\frac{a\cdot b}{\norm a\norm b}$ (scale-invariant), $\nu_1$ matches a
\emph{centered} spectrum at $\cos\approx0.92$, but only $0.67$ against the raw spectrum and $0.03$
against the mean: living in the centered space, $\nu_1$ is $\approx$orthogonal to the mean and is a
\emph{difference-spectrum} (Figure~\ref{fig:lead}), echoing Question~1's dominant-eigenvector steady state.

\begin{figure}[h]\centering
\includegraphics[width=0.74\linewidth]{q3d_v1_vs_row.png}
\caption{$\nu_1$ (green) overlaid with the best-matching centered spectrum ($\cos\approx0.92$); a
difference-spectrum on two molecular bands.}\label{fig:lead}
\end{figure}

\section{Part (e): clustering in covariance space}\label{sec:clusters}
With $V_4=[\nu_1\,\nu_2\,\nu_3\,\nu_4]$, the compressed coordinates $Z=\Xt V_4\in\R^{\Nt\times4}$ give
each vial four scores; the $\binom42=6$ pairwise scatters (Figure~\ref{fig:scatter}) show many tight,
separated groups. Counting by leader clustering (a numpy-only threshold method) gives $25$ clusters,
stable for radii in $[6,12]\times$ the median nearest-neighbour spacing --- a property of the data, not
the threshold. Thus the $924$ vials form $\approx25$ concoctions, below the ``fewer than $100$'' of
Part (h). Although $\nu_1,\dots,\nu_4$ live in feature space, projecting samples onto them exposes
sample-space structure, since similar spectra get similar scores. Part (f) repeats this with the Gram
matrix $G=\tfrac1{\Nt}\Xt\Xt\T$ (same nonzero eigenvalues as $C$, Question 2d), which should recover the
same groups.

\begin{figure}[h]\centering
\includegraphics[width=0.82\linewidth]{q3e_covariance_scatter.png}
\caption{The six pairwise scatters of the top-4 PC scores; each point is one vial. The data resolves
into $\approx25$ clusters.}\label{fig:scatter}
\end{figure}

\end{document}
